\section{Hypothesis and related work}\label{sec:related}

\paragraph{Betting markets underprice favourites.}
The oldest regularity in betting markets is that long shots are overpriced and favourites underpriced.
\citet{griffith1949} found it in American horse racing, and \citet{thaler1988} review the many markets where it has
since been confirmed: per dollar staked, bets on favourites lose less than bets on long shots, so a bettor who only
backs favourites does better than one who spreads money evenly. \citet{snowberg2010} use the prices of combination bets
to separate two explanations, a taste for risk and a misreading of small probabilities, and find that the misreading
explains the pattern better. Whatever its cause, the bias means that betting prices are \emph{underconfident}. They sit
too close to even odds, so the side priced as more likely wins more often than its price says.

\paragraph{Prediction markets show the same pattern.}
Prediction markets were designed to turn prices into forecasts. In theory a price need not equal the average belief of
traders \citep{manski2006}, but under broad conditions it comes close \citep{wolfers2006}, and in practice election
markets have long competed with polls: over five presidential elections, the Iowa Electronic Markets came closer to the
result than 964 polls in 74\% of comparisons \citep{berg2008}. Accuracy, however, is not calibration. A market can pick
the right winner and still put the wrong number on it. \citet{page2013} find that prediction-market prices are well
calibrated shortly before a market closes but biased further out, and they explain the bias by the favourite--longshot
pattern combined with the cost of tying up money until settlement. The exchanges that now dominate US election betting
show the same pattern. On Kalshi, cheap contracts win less often than their price requires and expensive ones earn small
positive returns, with different patterns for traders who post offers and traders who take them \citep{burgi2025}.
\citet{le2026} fits a recalibration slope to about 327{,}000 contracts, mainly on Kalshi, with a check on Polymarket. That slope is,
up to an intercept, the stretch we use. Political markets come out underconfident at nearly every time horizon, with a
mean slope of 1.64 on Kalshi and 1.31 on Polymarket. During the 2024 campaign, the same contracts also traded at
different prices on different exchanges \citep{clinton2025}, a reminder that these markets are far from frictionless.
The evidence therefore points one way: political prediction markets, like betting markets before them, are slightly
underconfident.

\paragraph{From a bias to a correction.}
A bias that is consistent can be corrected. If favourites win more often than their prices say, a forecaster should
push prices away from 50\%: a favourite priced at 80\% should be read as somewhat more than 80\%, and a long shot at 20\%
as somewhat less. The same idea runs through the forecasting literature. Averages of many forecasts tend to sit too
close to 50\%, and pushing them outwards makes them more accurate \citep{baron2014}. In election forecasting,
\citet{rothschild2009} corrects market prices for their known biases before comparing them with polls, and finds the
corrected prices more accurate than corrected polls, especially early in the campaign and in uncertain races. The stretch in this paper is the simplest such correction: one number that
pushes every price away from 50\% when it is above 1, and towards 50\% when it is below 1. Two practical questions follow
directly. First, does a stretch learned on past elections improve forecasts of the \emph{next} election, scored on
contests it has not seen? A correction that only fits the data it came from is of no use to a forecaster. Second, can
the bias be traded? If favourites are underpriced, buying them should earn more than they cost. The racetrack and Kalshi
evidence above suggests that any such edge is small, and fees and the gap between buying and selling prices may absorb
it.

\paragraph{Primaries look different.}
The studies above pool political markets, and our own data show why that matters. Of the 216 US contests in the training
period (152 on Kalshi, 64 on Polymarket), 101 are primaries. In general elections the familiar pattern holds: on 5 November 2024 the stretch was 1.97 on
Kalshi (95\% interval 1.27 to 3.77) and 1.83 on Polymarket (0.99 to no upper limit). In Kalshi's 78 primaries it was
0.32 (0.03 to 0.70 when whole election days are resampled): favourites there were not underpriced but heavily
overpriced, winning far less often than their prices said. Polymarket's 23 primaries were close to calibrated (0.88,
0.40 to 1.67). There are good reasons for primaries to differ. They are contests within one party, so the party label,
the strongest predictor in a general election, tells traders nothing. They often have several candidates, few or no
polls, low turnout and thin trading, and Kalshi lists many low-profile House primaries that receive little attention.
Pooled across both kinds of election, the two effects partly cancel, and Kalshi's markets even look overconfident
overall (0.58). A single correction for all election markets would then push one kind of election the wrong way. Primaries and
general elections therefore have to be studied, and corrected, separately.

\paragraph{Hypotheses.}
The training data leave one ambiguity. Most of its general elections are from 2024 and most of its primaries from 2026,
so the difference between them could be a difference between kinds of election or a change over time: perhaps the
underconfidence of 2024 had simply faded by 2026. The 2026 midterms, general elections held in 2026, separate the two
explanations. We examine three hypotheses, of which only the first was fixed before the vote.
\begin{itemize}
  \item \textbf{Main hypothesis (pre-registered).} General-election markets on each platform remain underconfident in
        the 2026 midterms, and are not less so than on 5 November 2024. If they do, primaries are traded differently
        from general elections. If instead they turn out calibrated, or overconfident, and below their 2024 value, the
        underconfidence faded over time. Section~\ref{sec:method} gives the exact test and the reading attached to each
        outcome.
  \item \textbf{Forecast correction (secondary).} Stretching midterm prices by a value fixed before the vote improves
        their accuracy. This was not part of the frozen test and is reported as an exploratory result.
  \item \textbf{Trading (exploratory).} Buying general-election favourites earns positive returns after realistic
        costs.
\end{itemize}
The 2026 midterms contain no primaries, so the primary result cannot yet be tested out of sample. The 2028 primary season
is where it can be, and we return to this in Section~\ref{sec:conclusion}.
